% Part: second-order-logic
% Chapter: sol-and-sets
% Section: power-of-continuum

\documentclass[../../../include/open-logic-section]{subfiles}

\begin{document}

\olfileid{sol}{set}{pow}

\olsection{सांतत्याचे आकारमान}

\begin{explain}
द्वितीय-क्रम तर्कशास्त्रात !!{domain} च्या उपसंचांवर
संख्यापन करता येते; पण !!{domain} च्या उपसंचांच्या
संचांवर करता येत नाही. ते थेट करण्यासाठी
\emph{तृतीय-क्रम} तर्कशास्त्र लागेल. उदाहरणार्थ,
संचाच्या घातसंचापासून मूळ संचाकडे !!{injective}
फलन नसते हे कँटरचे प्रमेय मांडायचे असल्यास आपण
असे म्हणण्याचा प्रयत्न करू शकतो: ``प्रत्येक
संच~$X$ आणि प्रत्येक संच~$P$ यांसाठी, $P$
हा~$X$ चा घातसंच असेल, तर
$\cardle{P}{X}$ नसते.'' पण $P$ हा~$X$ चा
घातसंच आहे असे म्हणण्यासाठी $P$ चे !!{element}s
नेमके~$X$ चे सर्व उपसंच आहेत, हे सूत्रबद्ध करावे
लागेल; उदाहरणार्थ,
$\lforall[Y][(P(Y) \liff Y \subseteq X)]$.
अडचण $P(Y)$ मध्ये आहे: ते द्वितीय-क्रम
तर्कशास्त्रातील !!a{formula} नाही, कारण
$P$ सारख्या एकस्थानी संबंधचलाचे निविष्ट फक्त
पदेच असू शकतात.

मात्र प्रांत पुरेसा मोठा असेल, तर संचांच्या संचांवरचे
संख्यापन \emph{अनुकरणाने} करता येते. कल्पना अशी
की द्विस्थानी संबंध~$R$ हा !!{domain} मधील
!!{element}s ना त्याच !!{domain} मधील
!!{element}s शी जोडतो. असा~$R$ दिल्यास,
एखाद्या~$x$ शी $R$-संबंधित सर्व !!{element}s
गोळा करता येतात: $\Setabs{y \in
  \Domain{M}}{R(x,y)}$ हा~$x$ ने
``सांकेतिक केलेला'' संच आहे. उलट,
$Z \subseteq \Pow{\Domain{M}}$ हा~$\Domain{M}$
च्या उपसंचांचा संग्रह असेल आणि~$\Domain{M}$
मध्ये~$Z$ मधील संचांइतके तरी !!{element}s
असतील, तर असा संबंध~$R \subseteq \Domain{M}^2$
असतो की प्रत्येक $Y \in Z$ हा
काही~$x$ ने~$R$ द्वारे सांकेतिक केला जातो.
\end{explain}

\begin{defn}
$R \subseteq \Domain{M}^2$ असेल, तर
$\Setabs{y \in \Domain{M}}{R(x, y)}$ या संचाचे
\emph{$x$ ने $R$-सांकेतीकरण} होते.
\end{defn}

!!a{element}~$x \in \Domain{M}$ हा
$Z \subseteq \Domain{M}$ संचाचे $R$-सांकेतीकरण
करत असेल, तर संच $Y \subseteq \Domain{M}$ हा
संचांच्या संचाचे सांकेतीकरण करतो: ते म्हणजे~$Y$
मधील !!{element}s नी सांकेतिक केलेले संच.
म्हणून संच $Y$ हा $\Pow{X}$ चे
$R$-सांकेतीकरण करू शकतो. असे होईल, जर
आणि फक्त जर प्रत्येक $Z \subseteq X$ साठी
काही $x \in Y$ हे~$Z$ चे $R$-सांकेतीकरण
करत असेल आणि प्रत्येक $x \in Y$ हा
काही~$Z \subseteq X$ चे $R$-सांकेतीकरण करत असेल.

\begin{prop}
!!{formula}
  \[
  \fn{Codes}(x, R, Z) \ident \lforall[y][(Z(y) \liff
    R(x, y))]
  \]
हे $s(x)$ हा $s(R)$ द्वारे $s(Z)$ चे
सांकेतीकरण करतो, असे व्यक्त करते. पुढील
!!{formula}
\begin{multline*}
  \fn{Pow}(Y, R, X) \ident \\
  \lforall[Z][(Z \subseteq X \lif \lexists[x][(Y(x) \land
    \fn{Codes}(x, R, Z))])] \land {} \\ \lforall[x][(Y(x) \lif
  \lforall[Z][(\fn{Codes}(x, R, Z) \lif Z \subseteq X)]]
\end{multline*}
हे $s(Y)$ हा $s(R)$ द्वारे~$s(X)$ च्या
घातसंचाचे सांकेतीकरण करतो, असे व्यक्त करते;
म्हणजे $s(Y)$ चे घटक $s(R)$ द्वारे
नेमके $s(X)$ चे सर्व उपसंच सांकेतिक करतात.
\end{prop}

\begin{explain}
या युक्तीने उपसंचांवर थेट संख्यापन करण्याऐवजी
त्यांच्या संकेतांवर संख्यापन करून घातसंचाविषयी
विधाने व्यक्त करता येतात. उदाहरणार्थ,
आता कँटरचे प्रमेय असे मांडता येते की~$X$ चा
घातसंच सांकेतिक करणाऱ्या कोणत्याही संबंधाच्या
प्रांतापासून मूळ~$X$ कडे !!{injective} फलन नसते.
\end{explain}

\begin{prop}
पुढील वाक्य
\begin{multline*}
  \lforall[X][\lforall[Y][\lforall[R][(\fn{Pow}(Y, R, X) \lif \\
      \lnot \lexists[u][(\lforall[x][\lforall[y][(\eq[u(x)][u(y)] \lif
            \eq[x][y])]] \land {}\\
        \lforall[x][(Y(x) \lif X(u(x)))])])]]]
\end{multline*}
वैध आहे.
\end{prop}

\begin{explain}
!!a{denumerable} संचाचा घातसंच !!{nonenumerable}
असतो; म्हणून त्याचा संचांक कोणत्याही
!!{denumerable} संचाच्या संचांकापेक्षा
मोठा असतो (तो संचांक $\aleph_0$ असतो).
$\Pow{\Nat}$ चे आकारमान ``सांतत्याचे आकारमान''
म्हणतात, कारण ते वास्तव संख्यारेषेवरील
बिंदूंच्या संचाच्या, म्हणजे~$\Real$ च्या,
आकारमानाइतके असते. एखाद्या
!!a{denumerable} संचाचा घातसंच सांकेतिक
करण्याइतका !!{domain} मोठा असेल, तर एखादा
संच सांतत्याच्या आकारमानाचा आहे हे असे
मांडता येते: तो~$\aleph_0$ आकारमानाच्या संच~$X$
चा घातसंच सांकेतिक करणाऱ्या कोणत्याही
संच~$Y$ शी तुल्यबल आहे. (प्रांत पुरेसा मोठा
नसेल, म्हणजे त्यात~$\Real$ शी तुल्यबल
उपसंचच नसेल, तर~$\Pow{X}$ चे सांकेतीकरण
करणारा संबंधही असू शकत नाही.)
\end{explain}

\begin{prop}
$\cardle{\Real}{\Domain{M}}$ असेल, तर पुढील
!!{formula}
\begin{multline*}
\fn{Cont}(Y) \ident
\lexists[X][\lexists[R][((\fn{Aleph}_0(X) \land
      \fn{Pow}(Y, R, X)) \land \\
      \lforall[x][\lforall[y][((Y(x) \land {}
      Y(y) \land \lforall[z][R(x,z) \liff R(y,z)]) \lif \eq[x][y])]])]]
\end{multline*}
हे $\cardeq{s(Y)}{\Real}$ असे व्यक्त करते.
\end{prop}

\begin{proof}
  $\fn{Pow}(Y, R, X)$ यानुसार $s(Y)$ हा $s(R)$
  द्वारे~$s(X)$ चा घातसंच सांकेतिक करतो;
  $\fn{Aleph}_0(X)$ नुसार हा संच गणनीय अनंत असतो.
  म्हणून $s(Y)$ चे आकारमान किमान सांतत्याच्या
  आकारमानाइतके असते; ते अधिकही असू शकते
  (जर $s(Y)$ चे अनेक घटक~$X$ चा तोच उपसंच
  सांकेतिक करत असतील). शेवटचे संयुक्तपद
  हे टाळते: $s(R)$ द्वारे $s(Y)$ चे घटक आणि
  $s(X)$ चे उपसंच यांच्यातील संगती
  !!{injective} असावी, अशी त्याची अट आहे.
  संपादकीय नोंद: आधीच्या विभागातील गणनीयतेचे
  सूत्र उपसंच हाच संपूर्ण प्रांत असावा असे
  भाग पाडते. मग त्याच प्रांतात त्या उपसंचाचा
  घातसंच सांकेतिक करण्याची येथील अट
  कँटरच्या प्रमेयाशी विसंगत होते. म्हणून
  सांतत्याचे हे सूत्र कोणत्याही रचनेत
  खरे ठरत नाही.
\end{proof}

\begin{prop}
$\cardeq{\Domain{M}}{\Real}$ जर आणि फक्त जर
\begin{multline*}
  \Sat{M}{\lexists[X][\lexists[Y][\lexists[R][(\fn{Aleph}_0(X) \land \fn{Pow}(Y, R, X) \land \\
          \lexists[u][(\lforall[x][\lforall[y][(\eq[u(x)][u(y)] \lif \eq[x][y])]] \land {} \\\lforall[y][(Y(y) \lif \lexists[x][\eq[y][u(x)]])])])]]]}.
        \end{multline*}
  संपादकीय नोंद: मागील गणनीयतेचे सूत्र
  उपसंचाला संपूर्ण प्रांत ठरवते, तर येथील
  घातसंचाचे सूत्र त्या प्रांताचा घातसंच
  त्याच प्रांतात सांकेतिक करायला सांगते.
  कँटरच्या प्रमेयानुसार हे अशक्य आहे.
  म्हणून उजवीकडील सूत्र कोणत्याही
  रचनेत सत्य होत नाही आणि ही समतुल्यता
  या स्वरूपात चुकीची आहे.
  \end{prop}

\begin{explain}
सांतत्य परिकल्पनेनुसार सांतत्याचे आकारमान
हा पहिला !!{nonenumerable} संचांक असतो;
म्हणजे $\Pow{\Nat}$ चे आकारमान~$\aleph_1$ असते.
\end{explain}

\begin{prop}
सांतत्य परिकल्पना खरी असेल, जर आणि फक्त जर
\[\fn{CH} \ident
\lforall[X][(\fn{Aleph}_1(X) \liff \fn{Cont}(X))]\]
वैध असेल, असा मूळ मजकुराचा दावा आहे.
संपादकीय नोंद: येथे वापरलेले अलेफ-एक
आणि सांतत्याची आकारमाने ओळखणारे आधीचे
सूत्रदावे सदोष असल्याने ही समतुल्यता
या स्वरूपात पडताळलेली नाही.
\end{prop}

हे लक्षात घ्या की $\lnot \fn{CH}$ वैध असेल,
जर आणि फक्त जर सांतत्य परिकल्पना खोटी असेल,
असे होत नाही. !!a{enumerable} प्रांतात
आकारमान~$\aleph_1$ असलेले उपसंच नसतात,
आणि सांतत्याच्या आकारमानाचे उपसंचही नसतात;
म्हणून अशा !!a{enumerable} प्रांतात
$\fn{CH}$ नेहमी खरे असते, असे मूळ
मजकुराचे स्पष्टीकरण आहे. मात्र सांतत्य
परिकल्पना खोटी असेल, जर आणि फक्त जर
वैध असेल असे वेगळे वाक्य देता येते,
असा पुढील दावा आहे:

\begin{prop}
सांतत्य परिकल्पना खोटी असेल, जर आणि
फक्त जर \[\fn{NCH} \ident
\lforall[X][(\fn{Cont}(X) \lif \lexists[Y][(Y \subseteq X \land \lnot
    \fn{Count}(Y) \land \lnot \cardeq{X}{Y})])]\]
वैध असेल, असा मूळ मजकुराचा दावा आहे.
संपादकीय नोंद: या दाव्यातही आधीच्या
संचांक-सूत्रांतील दोषांचा परिणाम तपासणे
आवश्यक आहे.
\end{prop}

\end{document}
